\section{Choice of Parameters}
\label{sec:param}

We choose the parameters of the construction so that the hypotheses of the
statements of the earlier sections that Section~\ref{sec:bootstrap} applies, in
particular the lemmas of Section~\ref{sec:relay}, hold at once. These statements
give their hypotheses explicitly, as inequalities between the parameters and the
constants of the earlier sections and of Part~I, and Section~\ref{sec:bootstrap}
assumes the choices made here when it applies them.

Recall that $\mu$ fixes the ratio of the amplitudes of consecutive rings and,
with $\sbar$, the ratios $s_k$ of their radii in Equation~\eqref{eq:datum}, and
that $\kseed$ is the seed size of Definition~\ref{def:approx}. The numbers
$\delta_0$ and $\Tpin$ are those of the clock of Definition~\ref{def:clock}. The choices are made in the
order recorded in Table~\ref{tab:param}, and Lemma~\ref{lem:nonempty} shows that
they can be made. Throughout this section we assume
\textup{(N1)}--\textup{(N3)}.

\begin{definition}
\label{def:P0-P4}
\emergencystretch=3em
The \textbf{structural list} consists of the following choices; the objects of
(SA1), (SA2), and the quantities that Definition~\ref{def:standing} derives from
them, are defined there; the list only fixes the order in which they are chosen.
\begin{description}[font=\normalfont]
\item[\textup{(P0.1)}] $m=3$; $\Phi$, $a$, $\Gone$ of Definition~I.\ref{I-def:ring}; $z$, $\sigma$,
  $\sigb$ of Definition~I.\ref{I-def:eigendata}.
\item[\textup{(P0.2)}] An integer $s$ and a number $\sigma'$ as in (SA1); then
  $c_\gamma$ of Definition~\ref{def:standing}.
\item[\textup{(P0.3)}] $p=m$; the integers $N_G$, $M$, the constants
  $\dH,C_2,C_3,\etmax$ and $\sigma_{\max}$, and the functions
  $\kappa_0(\cdot)$, $\bar t_1(\cdot)$, $C(\cdot)$ of Theorem~\ref{thm:active};
  $\kL$ and $\Cseed$ of Lemma~\ref{lem:seed}; $\|\ellst\|$, $\Cls$ and $\Cp$ of
  Lemma~I.\ref{I-lem:uniform}.
\item[\textup{(P0.4)}] The integers $s_h$, $N_Q$, $M_h$ and the constants
  $\majst,C_Q$ of (SA2).
\item[\textup{(P0.5)}] At order $M_h$: the constants $C_{L1}$, $C_M$, $C_O$,
  $\Cin$ of (SA2), and $C_\eta$ of Lemma~\ref{lem:field-bound}. Further $\Mvort$,
  $\Cstar$ of Lemma~\ref{lem:speed}; $C_G$ of Lemma~\ref{lem:inner-circ}; $\Cmom$
  of Proposition~\ref{prop:confine}.
\item[\textup{(P0.6)}] $c_N$, $C_N$, $\cout$ of Lemma~I.\ref{I-lem:uniform};
  $C_{\mathrm{qr}}$, $C_{\mathrm{id}}$, $C_{\mathrm{od}}$
  and the function $C_{\mathrm{tr}}(\cdot)$ of Lemma~\ref{lem:kicktrue}.
\item[\textup{(P1)}] Choose $\delta_0\in(0,\dH]$ sufficiently small that
  $\frac{16}{15}\cdot\frac{2\pi C_3\delta_0^2}{\sigma}\le\frac14\cdot\frac{15}{16}\cdot\frac{\cout\delta_0}{20|z|}$
  and $\Cp C_2\delta_0\le\frac{\cout}{16\pi}$.
\item[\textup{(P2)}] $\mu\in(1,\mu_{\max}]$, with
  $\mu_{\max}=1+\min\big(c'\delta_0,\frac1{16}\big)$ and
  a fixed $c'>0$ sufficiently small that
  $c'\le\frac{15}{16}\cdot\frac{\cout}{10|z|\,\Cmom\,\Ip{m}{\sigma/4}}$; then
  $\vartheta_\mu$, $\gamma_\mu$, $\gamma'_\mu$, $\CDuh$ of Definition~\ref{def:standing},
  and $\KO$ of Lemma~\ref{lem:O}.
\item[\textup{(P3)}] $\Tpin>0$ with $\rT(\Tpin)\le\frac1{64}$, with the function
  $\rT$ of Definition~\ref{def:standing}. Then $C_D(\Tpin)$
  (Proposition~\ref{prop:relay}); $c_K$, $K_1$, $\Abar$, $\Klo$, $\Kup$, $\Chand$
  (Definition~\ref{def:K1});
  $C_F\ge\Cin\,\mu\,(2+2\Cstar\Tpin)^m(1+2\Cstar)^m/\Klo$; $\Cdeep$
  (Lemma~\ref{lem:kicktrue}); $\Tz$ (Lemma~\ref{lem:time}) and
  choose $C_\tau\ge\max\{3,\frac{2m+5}{\sigma'},\frac{3(m+2)}{\sigma'}\}$.
\item[\textup{(P4)}] Choose $C_0\ge\max\{1,\,2\Cseed,\,\Chand,\,C_F\}$. Then $\kappa_0(C_0)$,
  $\bar t_1(C_0)$ and $C(C_0)$ (Theorem~\ref{thm:active}); $C_{\mathrm{tr}}(C_0)$
  (Lemma~\ref{lem:kicktrue}); $C_\Xi(C_0,K_1)$ (Proposition~\ref{prop:relay});
  $\Cop(C_0,K_1)$ (Lemma~\ref{lem:Oprime}).
\end{description}
\end{definition}

The number $C_0$ of (P4) is the constant of hypotheses (Init) and (Field) of
Theorem~\ref{thm:active}, which the proof of Proposition~\ref{prop:bootstrap}
applies to every active ring. That proof bounds the initial perturbation of the
seeded ring by $2\Cseed\ksh_n$, that of ring $k$ at its hand-off by
$\Chand\ksh_k$, and the inner field of the active ring $k$ by
$C_F\ksh_k\jap{u}^m$. Since (Init) and (Field) ask for bounds $C_0\kappa$ and
$C_0\kappa\jap{u}^p$, with $p=m$, the number $C_0$ must be at least $2\Cseed$,
$\Chand$ and $C_F$, and at least $1$ as Theorem~\ref{thm:active} requires. None of
$\Chand$, $C_F$, $\Cseed$ depends on $C_0$, so the choice of $C_0$ is not
circular. Likewise $\Tpin$ depends only on $\sigma$ and $m$. It is chosen before
$C_F$, whose formula in (P3) contains it. This order is not circular either: the
condition on $\Tpin$ in (P3) is the hypothesis $\rT(\Tpin)\le\frac1{64}$ of
Lemmas~\ref{lem:drift} and~\ref{lem:kicktrue}, and the drift bound of
Lemma~\ref{lem:drift} does not involve $C_F$.

\begin{definition}
\label{def:P5}
The \textbf{$\kappa$-list (P5)}: choose $\kseed>0$ so small that, for every
$\kappa\in(0,2\kseed]$, the following hold.
\begin{description}[font=\normalfont]
\item[\textup{($\kappa$1)}] $\kappa\le\delta_04^{-\mu/(\mu-1)}$.
\item[\textup{($\kappa$2)}] \begin{multline*}
  \frac{2\pi C_3\kappa^{1/2}\mu}{\frac34\sigma(\mu-1)}+C_{\mathrm{tr}}(C_0)\Big(\frac\kappa{\delta_0}\Big)^{\gamma_\mu}\log^{m+1}\frac{\delta_0}\kappa\\
  +\big[\Cdeep+C_{\mathrm{id}}\rT(\Tpin)\Ip{m}{\tfrac\sigma8}\big]\Big(\frac\kappa{\delta_0}\Big)^{\vartheta_\mu}\log^{m}\frac{\delta_0}\kappa
  \le\frac12\cdot\frac{15}{16}\cdot\frac{\cout}{20|z|(\mu-1)}.
  \end{multline*}
\item[\textup{($\kappa$3)}] $\Cpre(\kappa/\delta_0)^{\vartheta_\mu}\le\frac{15}{16}\cdot\frac{\cout}{20|z|(\mu-1)}$,
  where $\Cpre=C_{\mathrm{qr}}\big(\frac{4\mu C_{L1}C_Q}{3\sigma(\mu-1)}+1\big)K_1$.
\item[\textup{($\kappa$4)}] $\kappa\le\kappa_0(C_0)$.
\item[\textup{($\kappa$5)}] For every $G\in[0,\Gone]$, with $\hat T=\sigma(G)^{-1}\log(\delta_0/\kappa)$
  and $\rhohat$ the transient of Theorem~\ref{thm:active} for $(C_0,G)$:
  $\Cp C_2\big(\rhohat(\hat T)+\kappa^{1/2}\big)\le\frac{\cout}{16\pi}$.
\item[\textup{($\kappa$6)}] $\kappa\le\kL$.
\item[\textup{($\kappa$7)}] $C_\Xi(C_0,K_1)\big(1+\frac1{\sigb(\mu-1)}\big)(\kappa/\delta_0)^{\gamma'_\mu}\langle\log(\delta_0/\kappa)\rangle^{m+1}\le\frac1{\sigb(\mu-1)}$.
\end{description}
\end{definition}

In (P5), $\kappa$ stands for an amplitude $\ksh_k$ with which a ring starts its
active phase and, in ($\kappa$1) and ($\kappa$6), also for $2\kseed$.
Theorem~\ref{thm:active} is applied to ring $k$ with $\kappa=\ksh_k$, and
Lemma~\ref{lem:kicktrue}, at level $k$, with $\kappa=\ksh_{k+1}$.
Conditions ($\kappa$1)--($\kappa$3) are hypotheses (c)--(e) of
Lemma~\ref{lem:kicktrue}, ($\kappa$4) is hypothesis (Small) of
Theorem~\ref{thm:active}, ($\kappa$5) is used in the proof of
Proposition~\ref{prop:reduction}, and ($\kappa$6) makes the seed of
Definition~\ref{def:approx} admissible. Every amplitude $\ksh_k$ lies in
$(0,2\kseed]$ (Steps~0 and~3 of the proof of Proposition~\ref{prop:bootstrap}),
so the conditions are imposed on that whole interval.

Condition ($\kappa$7) is imposed after $C_0$ and $K_1$ are fixed. Here $K_1$
(Definition~\ref{def:K1}) is an explicit function of quantities fixed at
(P0)--(P3); it does not involve $C_0$, $\kseed$ or the deeper rings. Ring $j+1$
enters the hand-off of ring $j$ only through the quantity $\Xi_{j+1}$ of
Proposition~\ref{prop:relay}, and ($\kappa$7) at $\kappa=\ksh_{j+1}$ makes
$\Xi_{j+1}$ at most $\delta_0/(\sigb(\mu-1))$ at every level (Step~3 of the
proof of Proposition~\ref{prop:bootstrap}). So the bound on ring $j$ at its
hand-off holds with the same constants $K_1$ and $\Abar$ at every level.

The number $\sbar$, which fixes the ratios $s_k$ of consecutive radii, is chosen
last, because the conditions (B1)--(B11) on it involve constants fixed at all the
earlier steps; each of them holds for every
sufficiently small $\sbar$ (proof of Lemma~\ref{lem:nonempty}).

\begin{definition}
\label{def:P6}
The \textbf{$\sbar$-list (P6)}: choose $\sbar>0$ such that the following hold,
where $s_k=\sbar(2\mu)^{-k}$, $\eta_{b,1}=0$ and $\eta_{b,k}$, $\eta_b$ are as in
Equation~\eqref{eq:etab}.
\begin{description}[font=\normalfont]
\item[\textup{(B1)}] $\sbar\le a/16$ and $\sbar^{m-2}\le\frac12$.
\item[\textup{(B2)}] $\mu^4\sbar^2\le\frac18$.
\item[\textup{(B3)}] $C_G\mu\sbar^2\le\Gone$.
\item[\textup{(B4)}] $\sup_{0\le G\le C_G\mu\sbar^2}|\lamst(G)-z|\le\sigma(\mu-1)/(8\mu)$
  (this is (Det)).
\item[\textup{(B5)}] For every $k\ge1$: $\Lconf_ks_k\le a/8$ and
  $\frac83(\frac2a)^m(\Lconf_k)^{m-1}s_k^{m-2}\le\mu\Cstar$, where
  \[
  \Lconf_k=2+2\Cstar\mu\Big[\Tpin+1+\frac1{\sigma'}\Big((m+2)\log\frac1{s_k}+\log_+\frac{\delta_0}{\Klo}\Big)\Big].
  \]
\item[\textup{(B6)}] $\Kup\sbar^{m+2}\le\kseed$.
\item[\textup{(B7)}] $16C_OC_Q\CDuh\KO\sbar^{2m}\le1$, $\KO\sbar^{2m}\le\frac14$,
  $\KO\mu\sbar^{m+2}\le\frac12$, and
  \begin{equation}
  \label{eq:E22}
  \KO\sbar^{2m}\Big[\Tz+2\log\frac1{\sbar}+\frac1{\mu-1}\Big]\le\frac{\mu-1}{4C_OC_QC_\tau}\min\Big\{\frac12,\frac1{4\Cls K_1}\Big\}.
  \end{equation}
\item[\textup{(B8)}] $\eta_b\le\frac12\etmax$; for every $k\ge2$,
  \[
  \frac{\Cls\eta_{b,k}}{\sigma'}\Big[(m+2)\log\frac1{s_k}+\log_+\frac{\delta_0}{\Klo}\Big]\le\frac1{32};
  \]
  and $\Cls\eta_b\sigma'^{-1}\log(2\delta_0/\kseed)\le\frac1{32}$.
\item[\textup{(B9)}] $2\Abar\mu\sbar^{m+2}\le\majst$ and $C_\eta(\mu\sbar^2+\sbar^{m-2})\le1$.
\item[\textup{(B10)}] $4C_M\big(\frac{2C_Q}{\sigb(\mu-1)}+1\big)\mu\sbar^{m+2}\le\frac12$.
\item[\textup{(B11)}] $\|\ellst\|C_QK_1\Abar\mu\sbar^{m+2}\le\frac12$, and
  \[
  C_{\mathrm{od}}\eta_b\delta_0\frac{\mu^2}{\sigma^2(\mu-1)^2}+2\pi C_3\eta_b\frac{4\mu\delta_0}{3\sigma(\mu-1)}
  \le\frac14\cdot\frac{15}{16}\cdot\frac{\cout\delta_0}{20|z|(\mu-1)} .
  \]
\end{description}
\end{definition}

The second part of (B1), and (B2), hold automatically (see the remark after
Equation~\eqref{eq:sk-small}); they are listed because Proposition~\ref{prop:confine} and
Lemma~\ref{lem:field-bound} state them as hypotheses.

\begin{table}[!tb]
\centering
\small
\begin{tabular}{@{}p{0.40\textwidth}lp{0.42\textwidth}@{}}
\hline
\raggedright Quantity & Fixed at & \raggedright Determined by \tabularnewline
\hline
\raggedright $s$, $\sigma'$, $c_\gamma$ & (P0.2) & \raggedright $\sigb$, $\sigma$ ($s\ge3$ free) \tabularnewline
\raggedright $N_G$, $M$, $\dH$, $C_2$, $C_3$, $\etmax$, $\sigma_{\max}$, $\kappa_0(\cdot)$, $\bar t_1(\cdot)$, $C(\cdot)$ & (P0.3) & \raggedright $\Phi$, $\Gone$, $s$, $\sigma'$ \tabularnewline
\raggedright $\kL$; $\Cseed$; $\|\ellst\|$, $\Cls$, $\Cp$ & (P0.3) & \raggedright $\Phi$; $\Phi$, $\Gone$, $s$; $\Phi$, $\Gone$ \tabularnewline
\raggedright $s_h$, $N_Q$, $M_h$, $\majst$, $C_Q$ & (P0.4) & \raggedright $M$, $\Lam$, $\sigb$ \tabularnewline
\raggedright $C_{L1}$, $C_M$, $C_O$, $\Cin$, $C_\eta$ & (P0.5) & \raggedright $a$, $m$, $M_h$, $\|\Phi\|_{L^1}$ \tabularnewline
\raggedright $\Mvort$, $\Cstar$, $C_G$, $\Cmom$ & (P0.5) & \raggedright absolute; $\|\Phi\|_{L^1}$ \tabularnewline
\raggedright $c_N$, $C_N$, $\cout$, $C_{\mathrm{qr}}$, $C_{\mathrm{id}}$, $C_{\mathrm{od}}$, $C_{\mathrm{tr}}(\cdot)$ & (P0.6) & \raggedright $\Phi$, $\Gone$, $m$; $C_3$, $\sigma$, $C(\cdot)$ \tabularnewline
\raggedright $\delta_0$ & (P1) & \raggedright $\dH$, $C_2$, $C_3$, $\Cp$, $\cout$, $\sigma$, $|z|$ \tabularnewline
\raggedright $c'$, $\mu_{\max}$, $\mu$; $\vartheta_\mu$, $\gamma_\mu$, $\gamma'_\mu$, $\CDuh$, $\KO$ & (P2) & \raggedright $\delta_0$, $\cout$, $|z|$, $\Cmom$, $\sigma$; $c_\gamma$, $\sigb$, $C_M$, $C_Q$ \tabularnewline
\raggedright $\Tpin$ & (P3) & \raggedright $\sigma$, $m$ \tabularnewline
\raggedright $C_D(\Tpin)$, $c_K$, $K_1$, $\Abar$, $\Klo$, $\Kup$, $\Chand$ & (P3) & \raggedright (P0)--(P2), $\Tpin$ \tabularnewline
\raggedright $C_F$, $\Cdeep$, $\Tz$, $C_\tau$ & (P3) & \raggedright (P0)--(P2), $\Tpin$, $\Klo$ \tabularnewline
\raggedright $C_0$ & (P4) & \raggedright $\Cseed$, $\Chand$, $C_F$ \tabularnewline
\raggedright $\kappa_0(C_0)$, $\bar t_1(C_0)$, $C(C_0)$, $C_{\mathrm{tr}}(C_0)$, $C_\Xi(C_0,K_1)$, $\Cop(C_0,K_1)$ & (P4) & \raggedright $C_0$, $K_1$ and (P0) \tabularnewline
\raggedright $\kseed$ & (P5) & \raggedright (P0)--(P4) \tabularnewline
\raggedright $\sbar$; $\eta_b$, $\eta_{b,k}$, $\theta_j$ & (P6) & \raggedright (P0)--(P5); $\Cop$, $\sbar$, $\mu$, $C_O$, $C_Q$, $\KO$ \tabularnewline
\hline
\end{tabular}
\caption{Where each constant of the construction is fixed, and the quantities
that determine it. Every entry of the last column is fixed at an earlier row or
at the same step. The depth $n$, the amplitude $A_1$ and the solution
$\omega^{(n)}$ appear in no row; a level enters only as the index of
$\eta_{b,k}$ and $\theta_j$.}
\label{tab:param}
\end{table}

The last inequality of (B7), Equation~\eqref{eq:E22}, is used only through the
following lemma. Recall that $\theta_j$ is the factor of Lemma~\ref{lem:O} that
bounds the field of the outer rings on ring $j$ by the Duhamel majorant
$\maj_j$, and that $\bar\tau^\sharp_j$ of Lemma~\ref{lem:quiet-time} bounds the
length of the quiet phase of ring $j$ in its own time. The lemma bounds their
product, which must be at most $\frac12$ for hypothesis \textup{(H2)} of
Proposition~\ref{prop:relay} and for the corresponding hypothesis of
Lemma~\ref{lem:Oprime}, and at most $1/(4\Cls K_1)$ for hypothesis
\textup{(ii)} of Lemma~\ref{lem:handoff}.

\begin{lemma}
\label{lem:theta-quiet}
Assume (SA) and \textup{(B7)}, and let $K_1$ be as in Definition~\ref{def:K1},
$\Tz$ and $C_\tau$ as in \textup{(P3)}, $\theta_j$ as in Lemma~\ref{lem:O} and
$\bar\tau^\sharp_j$ as in Lemma~\ref{lem:quiet-time}. Then $\theta_j\bar\tau^\sharp_j\le\min\{\frac12,\frac1{4\Cls K_1}\}$
for every level $1\le j<n$.
\end{lemma}

\begin{proof}
For $j=1$, $\theta_1=0$ by Lemma~\ref{lem:O}, and the claim holds. Let
$2\le j<n$. By Lemma~\ref{lem:O}, $\theta_j=4C_OC_Q\KO s_{j-1}^{2m}$.

We first bound $\bar\tau^\sharp_j$ by a multiple of $C_\tau$. By (SA3),
$1<\mu\le\frac{17}{16}$, so $\mu\log(2\mu)\le\frac{17}{16}\log\frac{17}8<1$.
The three terms of $\bar\tau^\sharp_j$ (Lemma~\ref{lem:quiet-time}) are those in
$\Tz$, in $\log(1/s_j)$ and in $(\mu-1)^{-2}$; we compare them with
those of $C_\tau[\Tz+\log(1/s_j)+(\mu-1)^{-1}]/(\mu-1)$. By the definition of
$C_\tau$ in \textup{(P3)} and the bound on $\mu\log(2\mu)$, the coefficients
satisfy $3\le C_\tau$,
$(2m+5)/\sigma'\le C_\tau$ and $(m+2)\mu\log(2\mu)/\sigma'\le(m+2)/\sigma'\le C_\tau$;
this gives $\bar\tau^\sharp_j\le C_\tau\big[\Tz+\log(1/s_j)+(\mu-1)^{-1}\big]/(\mu-1)$.
Next we pass from $s_j$ to $s_{j-1}$. Since
$s_{j-1}\le\sbar\le a/16<\frac1{2\mu}$, we have
$\log\frac1{s_j}=\log\frac1{s_{j-1}}+\log(2\mu)\le2\log\frac1{s_{j-1}}$, so
\[
\theta_j\bar\tau^\sharp_j\le\frac{4C_OC_QC_\tau}{\mu-1}\,\KO\,g(s_{j-1}),\qquad
g(x)=x^{2m}\Big[\Tz+2\log\frac1x+\frac1{\mu-1}\Big].
\]
It remains to replace $s_{j-1}$ by $\sbar$, since Equation~\eqref{eq:E22}
bounds $\KO g(\sbar)$; for this we show that $g$ is nondecreasing near $0$. For $0<x<1$,
$g'(x)=x^{2m-1}\big[2m\big(\Tz+2\log\frac1x+\frac1{\mu-1}\big)-2\big]$, which is
nonnegative when $\log(1/x)\ge1/(2m)$, because $\Tz>0$ and $\mu>1$. So $g$ is
nondecreasing on $(0,e^{-1/(2m)}]$. Since $s_{j-1}\le\sbar\le a/16<\frac1{16}<e^{-1/(2m)}$,
$g(s_{j-1})\le g(\sbar)$, and Equation~\eqref{eq:E22} gives
\[
\theta_j\bar\tau^\sharp_j\le\frac{4C_OC_QC_\tau}{\mu-1}\KO g(\sbar)\le\min\{\frac12,\frac1{4\Cls K_1}\}.\qedhere
\]
\end{proof}

\begin{lemma}
\label{lem:nonempty}
Assume \textup{(N1)}--\textup{(N3)}. The choices (P0)--(P6) can be made in this
order, with all quantities finite and positive (except $\eta_{b,1}=\theta_1=0$),
and each quantity is determined by quantities fixed at earlier steps or earlier in the same step, as in
Table~\ref{tab:param}. None of the choices depends on $n$, $A_1$ or
$\omega^{(n)}$. Finally, if (P0)--(P6) hold, then for every integer $n\ge2$ and
every $A_1>0$ the standing assumptions (SA) hold with these parameters, and
(Det) holds.
\end{lemma}

\begin{proof}
We take the steps in order. At each step the quantities of the earlier steps are
fixed, and we show that the conditions of the step can be met.

\emph{(P0)--(P4).} The constants of (P0) exist and are finite by the results
cited there; in particular, by the formulas of (SA2) and $M\ge s+4$
(Theorem~\ref{thm:active}), $s_h=M-1\ge s+3\ge3$ and $N_Q\sigb\ge\Lam(s_h)$, as
Theorem~\ref{thm:quiet} requires. The interval $\big(\frac9{200},\min(\sigb,\frac34\sigma)\big)$ for $\sigma'$
in (SA1) is nonempty: by Proposition~I.\ref{I-prop:U}\textup{(a)}, $\sigb$ is attained and
$\sigb>\frac{211}{2500}>\frac9{200}$; and
$\frac34\sigma\ge\frac34\sigb>\frac{633}{10000}>\frac9{200}$. Moreover
$c_\gamma>0$, because $\sigma'<\frac34\sigma<\frac{15}{16}\sigma$. The conditions in (P1) hold for sufficiently small $\delta_0>0$ because
$\sigma\ge\sigb>0$, $\cout>0$ and the other constants have already been fixed.
In (P2), $c'>0$ because $\Cmom>0$ and $0<\Ip{m}{\sigma/4}<\infty$, so
$\mu_{\max}>1$. In (P3), $\Tpin$ exists because $\rT(T)=(1+T)^me^{-\sigma T/4}$
(Definition~\ref{def:standing}) tends to $0$ as $T\to\infty$, as $\sigma>0$. The formulas of (P3) and (P4) have positive finite denominators:
$c_K=3c_N/(32\mu)>0$, $\delta_0>0$, $\sigb>0$, $\mu-1>0$ and $\Klo>0$. So all the
constants of (P0)--(P4) are finite and positive, and $C_0\ge1$ and
$0<C_F\le C_0$, as Theorem~\ref{thm:active} and Lemma~\ref{lem:Oprime} require.
Each can be chosen in terms of the quantities in its row of Table~\ref{tab:param}, which are fixed at
earlier rows or earlier in the same step.

\emph{(P5).} Since (P5) is imposed for every $\kappa\in(0,2\kseed]$, it suffices
to find $x_0>0$ such that each condition holds for every $\kappa\in(0,x_0]$; then every $\kseed\in(0,x_0/2]$ is admissible. For ($\kappa$1), ($\kappa$4), ($\kappa$6) this holds
because $\delta_04^{-\mu/(\mu-1)}$, $\kappa_0(C_0)$, $\kL$ are positive. For
($\kappa$2), ($\kappa$3), ($\kappa$7) it suffices that each left side tends to
$0$ as $\kappa\to0$. The left sides are finite sums of terms $C\kappa^{1/2}$ and
$C(\kappa/\delta_0)^\gamma\langle\log(\delta_0/\kappa)\rangle^N$ with $C$ fixed at
(P0)--(P4), integers $N\ge0$, and
$\gamma\in\{\gamma_\mu,\vartheta_\mu,\gamma'_\mu\}$, positive because $\mu>1$ and
$c_\gamma>0$; each term tends to $0$. For ($\kappa$5) we need a bound uniform in
$G$. Let $\kappa<\delta_0$ and $G\in[0,\Gone]$, and put $c=\sigb-\sigma'>0$, where
$\sigma'<\sigb$ by (SA1). By Theorem~\ref{thm:active},
$t_1(C_0,G)\le\bar t_1(C_0)$, and $\sigma(G)\in[\sigb,\sigma_{\max}]$, so
$\hat T\ge\sigma_{\max}^{-1}\log(\delta_0/\kappa)$. If moreover
$\kappa\le\delta_0e^{-\sigma_{\max}\bar t_1(C_0)}$, then $\hat T\ge\bar t_1(C_0)$, so
$y=\hat T-t_1(C_0,G)\ge\hat T-\bar t_1(C_0)\ge0$. Hence $\hat T$ lies in the range
$u\ge t_1$ of the second formula for the transient in Theorem~\ref{thm:active},
and, as $\sigma(G)-\sigma'\ge c$,
\[
\rhohat(\hat T)=\jap{y}^pe^{-(\sigma(G)-\sigma')y}\le\jap{y}^pe^{-cy}\le K_ce^{-cy/2}
\le K_ce^{c\bar t_1(C_0)/2}\Big(\frac\kappa{\delta_0}\Big)^{c/(2\sigma_{\max})},
\]
where $K_c=\sup_{y\ge0}\jap{y}^pe^{-cy/2}<\infty$, and the last step uses
$y\ge\hat T-\bar t_1(C_0)$ and the lower bound on $\hat T$. The right side does not depend on
$G$ and tends to $0$ as $\kappa\to0$, and so does $\kappa^{1/2}$; so ($\kappa$5)
holds for all small $\kappa$, uniformly in $G\in[0,\Gone]$. We take for $x_0$
the least of the thresholds found for ($\kappa$1)--($\kappa$7).

\emph{(P6)}, with (P0)--(P5) fixed. It suffices that each condition holds for
every $\sbar$ in some interval $(0,\sbar_i]$, $\sbar_i>0$; then
every $\sbar\in(0,\min_i\sbar_i]$ is admissible. The conditions (B1)--(B3), (B6), (B7), (B9),
(B10), (B11), the first part of (B8) and its last part (with $\kseed$ fixed)
have left sides that are finite sums of terms $C\sbar^\alpha\log^N(1/\sbar)$ with
$\alpha>0$, $N\in\{0,1\}$ and $C$ fixed at (P0)--(P5); here we use
$\eta_b=\Cop\sbar^{2m}/(\sigb(\mu-1))$, and $m=3$, so that the exponent $m-2$
in (B1) and (B9) is positive. Each such term tends to $0$ as
$\sbar\to0$, while the right sides are fixed positive numbers. For (B4), the
supremum tends to $0$ as $\sbar\to0$, because $\lamst$ is continuous at $G=0$
(Proposition~I.\ref{I-prop:U}\textup{(a)}) and $\lamst(0)=z$; the right side
$\sigma(\mu-1)/(8\mu)$ is positive. The conditions (B5) and (B8) are imposed at
every level $k$, and we reduce each of them to one condition on $\sbar$. For
(B5), put
$\alpha_0=2+2\Cstar\mu\big[\Tpin+1+\sigma'^{-1}\log_+(\delta_0/\Klo)\big]$ and
$\beta_0=2\Cstar\mu(m+2)/\sigma'$, so that $\Lconf_k=\alpha_0+\beta_0\log(1/s_k)$.
The functions $x\mapsto x(\alpha_0+\beta_0\log(1/x))$ and
$x\mapsto x^{m-2}(\alpha_0+\beta_0\log(1/x))^{m-1}=x(\alpha_0+\beta_0\log(1/x))^2$
tend to $0$ as $x\to0$, so there is $\sbar_{\mathrm{B5}}>0$ such that both inequalities of
(B5) hold at level $k$ whenever $s_k\le\sbar_{\mathrm{B5}}$. Since $s_k\le\sbar$ for every
$k$, (B5) holds for every $\sbar\le\sbar_{\mathrm{B5}}$. For (B8) and $k\ge2$, we have
$s_k=s_{k-1}/(2\mu)$, and, by the formula for $\eta_{b,k}$ in
Equation~\eqref{eq:etab}, the $k$-th inequality reads $F(s_{k-1})\le\frac1{32}$
with
\[
F(x)=\frac{\Cls\Cop}{\sigma'\sigb(\mu-1)}\,x^{2m}\Big[(m+2)\Big(\log\frac1x+\log(2\mu)\Big)+\log_+\frac{\delta_0}{\Klo}\Big].
\]
Since $F(x)\to0$ as $x\to0$, there is $\sbar_{\mathrm{B8}}>0$ with $F\le\frac1{32}$ on
$(0,\sbar_{\mathrm{B8}}]$, and, since $s_{k-1}\le\sbar$ for $k\ge2$, the
inequalities hold for all $k\ge2$ once $\sbar\le\sbar_{\mathrm{B8}}$.

No formula or condition of (P0)--(P6) contains $n$, $A_1$ or $\omega^{(n)}$; the
levels in (B5) and (B8) range over all integers $k\ge1$ and $k\ge2$, not only
up to $n$.

\emph{(SA) and (Det).} Assume (P0)--(P6) and let $n\ge2$ and $A_1>0$. (SA1) holds by (P0.2) and (P0.3), and (SA2) by (P0.4) and (P0.5). For
(SA3): $\mu\in(1,\frac{17}{16}]$ because
$\mu_{\max}\le1+\frac1{16}$; $\sbar\in(0,a/16]$ and $C_G\mu\sbar^2\le\Gone$ by (B1) and (B3);
$\delta_0\in(0,\dH]$ by (P1); $\Tpin>0$ by (P3); and $\kseed\in(0,\kL]$ by
($\kappa$6) at $\kappa=2\kseed$. The objects of
Definitions~\ref{def:datum}--\ref{def:clock} are then defined with these
parameters: Definition~\ref{def:datum} needs $\mu\in(1,2]$ and $\sbar\in(0,a/16]$,
Definition~\ref{def:approx} needs $\kseed\in(0,\kL]$, and
Definition~\ref{def:clock} needs $\delta_0,\Tpin>0$ and $G_k\le\Gone$ for
$1\le k\le n$, which holds by Lemma~\ref{lem:inner-circ} and (B3). Hypothesis
(Det) is (B4).
\end{proof}
